\section{Review}
\label{section-algebraic}

\noindent
In this section we give a quick review of what has transpired above.

\medskip\noindent
Let $K/k$ be a field extension. Let $\alpha \in K$. Then we have the
following possibilities:
\begin{enumerate}
\item The element $\alpha$ is transcendental over $k$.
\item The element $\alpha$ is algebraic over $k$. Denote
$P(T) \in k[T]$ its {\it minimal polynomial}. This is a monic polynomial
$P(T) = T^d + a_1 T^{d - 1} + \ldots + a_d$ with coefficients in
$k$. It is irreducible and $P(\alpha) = 0$. These properties
uniquely determine $P$, and the integer $d$ is called the
{\it degree of $\alpha$ over $k$}. There are two subcases:
\begin{enumerate}
\item The polynomial $\text{d}P/\text{d}T$ is not identically zero.
This is equivalent to the condition that
$P(T) = \prod_{i = 1, \ldots, d} (T - \alpha_i)$ for
pairwise distinct elements $\alpha_1, \ldots, \alpha_d$
in the algebraic closure of $k$.
In this case we say that $\alpha$ is {\it separable} over $k$.
\item The $\text{d}P/\text{d}T$ is identically zero. In this case the
characteristic $p$ of $k$ is $ > 0$, and $P$ is actually a polynomial
in $T^p$. Clearly there exists a largest power $q = p^e$ such that $P$ is
a polynomial in $T^q$. Then the element $\alpha^q$ is separable over $k$.
\end{enumerate}
\end{enumerate}

\begin{definition}
\label{definition-separable-algebraic}
Algebraic field extensions.
\begin{enumerate}
\item A field extension $K/k$ is called {\it algebraic}
if every element of $K$ is algebraic over $k$.
\item An algebraic extension $k'/k$ is called {\it separable}
if every $\alpha \in k'$ is separable over $k$.
\item An algebraic
extension $k'/k$ is called {\it purely inseparable} if
the characteristic of $k$ is $p > 0$ and for every element
$\alpha \in k'$ there exists a power $q$ of $p$ such that
$\alpha^q \in k$.
\item An algebraic extension $k'/k$ is called {\it normal}
if for every $\alpha \in k'$ the minimal polynomial $P(T) \in k[T]$
of $\alpha$ over $k$ splits completely into linear factors over $k'$.
\item An algebraic extension $k'/k$ is called {\it Galois}
if it is separable and normal.
\end{enumerate}
\end{definition}

\noindent
The following lemma does not seem to fit anywhere else.

\begin{lemma}
\label{lemma-pth-root}
Let $K$ be a field of characteristic $p > 0$. Let $L/K$ be a separable
algebraic extension. Let $\alpha \in L$.
\begin{enumerate}
\item If the coefficients of the minimal polynomial of $\alpha$
over $K$ are $p$th powers in $K$ then $\alpha$ is a $p$th
power in $L$.
\item More generally, if $P \in K[T]$ is a polynomial such that (a) $\alpha$
is a root of $P$, (b) $P$ has pairwise distinct roots in an algebraic closure,
and (c) all coefficients of $P$ are $p$th powers, then $\alpha$ is a
$p$th power in $L$.
\end{enumerate}
\end{lemma}

\begin{proof}
It follows from the definitions that (2) implies (1). Assume $P$ is as in (2).
Write $P(T) = \sum\nolimits_{i = 0}^d a_i T^{d - i}$ and $a_i = b_i^p$.
The polynomial $Q(T) = \sum\nolimits_{i = 0}^d b_i T^{d - i}$ has distinct
roots in an algebraic closure as well, because the roots of $Q$
are the $p$th roots of the roots of $P$. If $\alpha$ is not a $p$th power,
then $T^p - \alpha$ is an irreducible polynomial over $L$
(Lemma \ref{lemma-take-pth-root}).
Moreover $Q$ and $T^p - \alpha$ have a root in common in
an algebraic closure $\overline{L}$.
Thus $Q$ and $T^p - \alpha$ are not relatively prime, which
implies $T^p - \alpha | Q$ in $L[T]$. This contradicts the fact that
the roots of $Q$ are pairwise distinct.
\end{proof}












